% Emacs, this is -*-latex-*-

\ex{Infer\^encia variacional de campo m\'edio I}
\label{ex:MFVI-I}
Seja $\ELBOx(q)$ o limite inferior da evid\^encia para a marginal $p(\x)$ de uma fdp/fmp conjunta $p(\x,\y)$,
\begin{equation}
    \ELBOx(q) = \E_{q(\y|\x)} \left[\log \frac{p(\x,\y)}{q(\y|\x)}\right].
\end{equation}
A infer\^encia variacional de campo m\'edio (*mean field*) assume que a distribui\'c\~ao variacional $q(\y|\x)$ se fatoriza completamente, isto \'e,
\begin{equation}
    q(\y | \x) = \prod_{i=1}^d q_i(y_i | \x),
\end{equation}
quando $\y$ \'e $d$-dimensional. Uma abordagem para aprender os $q_i$ para
cada dimens\~ao \'e atualizar um de cada vez enquanto mant\'em os outros fixos. N\'os
aqui derivamos as equa\'c\~oes de atualiza\'c\~ao correspondentes.



\begin{exenumerate}
    \item Mostre que o limite inferior da evid\^encia $\ELBOx(q)$ pode ser escrito como
    \begin{equation}
        \ELBOx(q) = \E_{q_1(y_1|\x)} \E_{q(\y_{\setminus 1}|\x)}\left[ \log p(\x,\y)\right] - \sum_{i=1}^d \E_{q_i(y_i|\x)} \left[\log q_i(y_i|\x)\right]
    \end{equation}
    onde $q(\y_{\setminus 1}|\x) = \prod_{i=2}^d q_i(y_i | \x)$ \'e a distribui\'c\~ao variacional sem $q_1(y_1|\x)$.
    
    \begin{solution}
        Isso segue diretamente da defini\'c\~ao do ELBO e da fatoriza\'c\~ao assumida de $q(\y|\x)$. Temos
        \begin{align}
            \ELBOx(q) &= \E_{q(\y|\x)} \log p(\x,\y) - \E_{q(\y|\x)} \log q(\y|\x) \\
            & = \E_{ \prod_{i=1}^d q_i(y_i | \x)}  \log p(\x,\y) - \E_{ \prod_{i=1}^d q_i(y_i | \x)} \sum_{i=1}^d \log q_i(y_i|\x) \\
            & = \E_{ \prod_{i=1}^d q_i(y_i | \x)}  \log p(\x,\y) - \sum_{i=1}^d  \E_{q_i(y_i | \x)} \log q_i(y_i|\x) \\
            & = \E_{q_1(y_1|\x)} \E_{ \prod_{i=2}^d q_i(y_i | \x)}  \log p(\x,\y) - \sum_{i=1}^d  \E_{q_i(y_i | \x)} \log q_i(y_i|\x)\\
            & = \E_{q_1(y_1|\x)} \E_{q(\y_{\setminus 1}|\x)}\left[ \log p(\x,\y)\right] - \sum_{i=1}^d \E_{q_i(y_i|\x)} \left[\log q_i(y_i|\x)\right]
        \end{align}
        Usamos aqui a linearidade da esperan\'ca. No caso de vari\'aveis aleat\'orias cont\'inuas, por exemplo, temos
        \begin{align}
            \E_{ \prod_{i=1}^d q_i(y_i | \x)} \sum_{i=1}^d \log q_i(y_i|\x) & = \int  q_1(y_1 | \x)\cdot \ldots\cdot q_d(y_d|\x)  \sum_{i=1}^d \log q_i(y_i|\x) d y_1\ldots d y_d\\
            & =  \sum_{i=1}^d \int q_1(y_1 | \x)\cdot \ldots \cdot q_d(y_d|\x) \log q_i(y_i|\x) d y_1\ldots d y_d\\
            & =  \sum_{i=1}^d \int q_i(y_i | \x)\log q_i(y_i|\x) d y_i \underbrace{\int \prod_{j\neq i} q_j(y_j|\x) d y_j}_{=1}\\
            & =  \sum_{i=1}^d E_{q_i(y_i|\x)} \log q_i(y_i|\x)
        \end{align}
        Para vari\'aveis aleat\'orias discretas, a integral \'e substitu\'ida por uma soma e leva ao mesmo resultado.
        
    \end{solution}
    
    \item Assuma que gostar\'iamos de atualizar $q_1(y_1 |\x)$ e que as
    marginais variacionais das outras dimens\~oes s\~ao mantidas fixas. Mostre que
    \begin{equation}
        \argmax_{q_1(y_1|\x)} \ELBOx(q) = \argmin_{q_1(y_1|\x)} \KL(q_1(y_1|\x) || \bar{p}(y_1|\x))
    \end{equation}
    com
    \begin{equation}
        \log \bar{p}(y_1|\x) = \E_{q(\y_{\setminus 1}|\x)}\left[ \log p(\x,\y)\right] + \text{const},
    \end{equation}
    onde $\text{const}$ refere-se a termos que n\~ao dependem de $y_1$. Isto \'e,
    \begin{equation}
        \bar{p}(y_1|\x) = \frac{1}{Z} \exp\left[ \E_{q(\y_{\setminus
                1}|\x)}\left[ \log p(\x,\y)\right] \right],
    \end{equation}
    onde $Z$ \'e a constante de normaliza\'c\~ao. Note que as vari\'aveis $y_2, \ldots, y_d$
    s\~ao marginalizadas devido \`a esperan\'ca em rela\'c\~ao a $q(\y_{\setminus
        1}|\x)$.
    
    \begin{solution}
        Come\'cando de 
        \begin{equation}
            \ELBOx(q) = \E_{q_1(y_1|\x)} \E_{q(\y_{\setminus 1}|\x)}\left[ \log p(\x,\y)\right] - \sum_{i=1}^d \E_{q_i(y_i|\x)} \left[\log q_i(y_i|\x)\right]
        \end{equation}
        descartamos os termos que n\~ao dependem de $q_1$. Obtemos ent\~ao
        \begin{align}
            J(q_1) & =  \E_{q_1(y_1|\x)} \E_{q(\y_{\setminus 1}|\x)}\left[ \log p(\x,\y)\right] -  \E_{q_1(y_1|\x)} \left[\log q_1(y_1|\x)\right]\\
            & =  \E_{q_1(y_1|\x)} \left[\log \bar{p}(y_1|\x) - \text{const} \right]- \E_{q_1(y_1|\x)} \left[\log q_1(y_1|\x)\right]\\
            & = \E_{q_1(y_1|\x)} \log \bar{p}(y_1|\x) - \E_{q_1(y_1|\x)} \left[\log q_1(y_1|\x)\right] - \text{const}\\
            & = \E_{q_1(y_1|\x)}\left[ \log \frac{\bar{p}(y_1|\x)}{ q_1(y_1|\x)} \right] - \text{const}\\
            & = -\KL( q_1(y_1|\x) || \bar{p}(y_1|\x) ) - \text{const}
        \end{align}
        Portanto
        \begin{equation}
            \argmax_{q_1(y_1|\x)} \ELBOx(q) = \argmin_{q_1(y_1|\x)} \KL(q_1(y_1|\x) || \bar{p}(y_1|\x))
        \end{equation}
        
    \end{solution}
    
    \item Conclua que, dados $q_i(y_i|\x)$, $i=2, \ldots, d$, o $q_1(y_1|\x)$ \'otimo \'e igual a $\bar{p}(y_1|\x)$.
    
    Isso ent\~ao leva a um esquema de atualiza\'c\~ao iterativo onde alternamos
    pelas diferentes dimens\~oes, cada vez atualizando a margimal variacional correspondente de acordo com:
    \begin{align}
        q_i(y_i|\x) &=  \bar{p}(y_i|\x), &  \bar{p}(y_i|\x) &= \frac{1}{Z} \exp\left[ \E_{q(\y_{\setminus i}|\x)}\left[ \log p(\x,\y)\right] \right]
    \end{align}
    onde $q(\y_{\setminus i}|\x) = \prod_{j \neq i} q(y_j | \x)$ \'e o produto de todas as marginais exceto a marginal $q_i(y_i|\x)$.
    
    \begin{solution}
        Isso segue imediatamente do fato de que a diverg\^encia KL \'e
        minimizada quando $q_1(y_1|\x) = \bar{p}(y_1|\x)$. Nota lateral: O
        regra de atualiza\'c\~ao iterativa pode ser considerada uma otimiza\'c\~ao por
        ascens\~ao de coordenadas no espa\'co de fun\'c\~oes, onde cada ``coordenada''
        corresponde a um $q_i(y_i|\x)$.
        
    \end{solution}
\end{exenumerate}

% --------------------------------------------------

\ex{Infer\^encia variacional de campo m\'edio II}
Assuma que as vari\'aveis aleat\'orias $y_1, y_2, x$ sejam geradas de acordo com o seguinte processo
\begin{align}
    y_1 &\sim \Gauss(y_1; 0, 1) &   y_2 &\sim \Gauss(y_2; 0, 1) \\
    n &\sim \Gauss(n; 0, 1)    &    x & = y_1+y_2+n & 
\end{align}
onde $y_1, y_2, n$ s\~ao estatisticamente independentes.

\begin{exenumerate}
    
    \item $y_1, y_2, x$ s\~ao conjuntamente Gaussianas. Determine sua m\'edia e sua
    matriz de covari\^ancia.
    
    \begin{solution}
        
        O valor esperado de $y_1$ e $y_2$ \'e zero. Pela linearidade da
        esperan\'ca, o valor esperado de $x$ \'e
        \begin{equation}
            \E(x) = \E(y_1) + \E(y_2) + \E(n) = 0
        \end{equation}
        A vari\^ancia de $y_1$ e $y_2$ \'e 1. Como $y_1, y_2, n$ s\~ao estatisticamente independentes,
        \begin{equation}
            \Var(x) = \Var(y_1) + \Var(y_2) + \Var(n) = 1 + 1 + 1 = 3.
        \end{equation}
        A covari\^ancia entre $y_1$ e $x$ \'e
        \begin{align}
            \text{cov}(y_1, x) &= \E( (y_1-\E(y_1))(x-\E(x))) = \E( y_1 x) \\
            &= \E( y_1(y_1+y_2+n) ) = \E(y_1^2) + \E(y_1 y_2) + \E(y_1 n)\\
            &= 1 +  \E(y_1)\E(y_2) + \E(y_1)\E(n)\\
            &= 1 + 0 + 0
        \end{align}
        onde usamos que $y_1$ e $x$ possuem m\'edia zero e as hip\'oteses de independ\^encia.
        
        A covari\^ancia entre $y_2$ e $x$ \'e computada da mesma forma e tamb\'em \'e igual a 1.
        
        Obtemos assim a matriz de covari\^ancia $\Sigmab$,
        \begin{equation}
            \Sigmab = \begin{pmatrix}
                1 & 0 & 1\\
                0 & 1 & 1\\
                1 & 1 & 3
            \end{pmatrix}
        \end{equation}
        
    \end{solution}
    
    \item A condicional $p(y_1, y_2 | x)$ \'e Gaussiana com m\'edia $\m$ e covari\^ancia $\C$,
    \begin{align}
        \m &= \frac{x}{3} \begin{pmatrix}
            1\\
            1
        \end{pmatrix}
        &
        \C & = \frac{1}{3} \begin{pmatrix}
            2 &-1 \\
            -1 & 2
        \end{pmatrix}
    \end{align}
    Como $x$ \'e a soma de tr\^es vari\'aveis aleat\'orias que possuem a
    mesma distribui\'c\~ao, faz sentido intuitivo que a m\'edia atribua
    $1/3$ do valor observado de $x$ para $y_1$ e $y_2$. Al\'em disso,
    $y_1$ e $y_2$ s\~ao negativamente correlacionadas, pois um aumento em $y_1$
    deve ser compensado com uma diminui\'c\~ao em $y_2$.
    
    Vamos agora aproximar a posterior $p(y_1, y_2 | x)$ com infer\^encia variacional
    de campo m\'edio. Determine a distribui\'c\~ao variacional \'otima usando o m\'etodo e os resultados do \exref{ex:MFVI-I}. Voc\^e pode usar que
    \begin{align}
        p(y_1, y_2, x) &= \Gauss\left( (y_1, y_2, x); \zerob, \Sigmab \right) &  \Sigmab &= \begin{pmatrix} 
            1 & 0 & 1\\
            0 & 1 & 1\\
            1 & 1 & 3
        \end{pmatrix}
        &
        \Sigmab^{-1} &= \begin{pmatrix} 
            2 & 1 & -1\\
            1 & 2 & -1\\
            -1 & -1 & 1
        \end{pmatrix}
    \end{align}
    
    \begin{solution}
        A hip\'otese de campo m\'edio significa que a distribui\'c\~ao variacional \'e assumida como fatoriz\'avel como
        \begin{equation}
            q(y_1, y_2 | x) = q_1(y_1 | x) q_2(y_2 | x)
        \end{equation}
        Do \exref{ex:MFVI-I}, as $q_1(y_1|x)$ e $q_2(y_2|x)$ \'otimas satisfazem
        \begin{align}
            q_1(y_1|x) &=  \bar{p}(y_1|x), &  \bar{p}(y_1|x) &= \frac{1}{Z} \exp\left[ \E_{q_2(y_2|x)}\left[ \log p(y_1, y_2, x)\right] \right]\\
            q_2(y_2|x) &=  \bar{p}(y_2|x), &  \bar{p}(y_2|x) &= \frac{1}{Z} \exp\left[ \E_{q_1(y_1|x)}\left[ \log p(y_1, y_2, x)\right] \right]
        \end{align}
        Note que estas s\~ao equa\'c\~oes acopladas: $q_2$ aparece na
        equa\'c\~ao para $q_1$ via $\bar{p}(y_1|x)$, e $q_1$ aparece na
        equa\'c\~ao para $q_2$ via $\bar{p}(y_2|x)$. Mas temos duas
        equa\'c\~oes para duas inc\'ognitas, que para o modelo conjunto Gaussiano
        $p(x, y_1, y_2)$ podem ser resolvidas em forma fechada.
        
        Dada a equa\'c\~ao fornecida para $p(y_1, y_2, x)$, temos que
        \begin{align}
            \log p(y_1, y_2, x) & = -\frac{1}{2} \begin{pmatrix}
                y_1\\
                y_2\\
                x
            \end{pmatrix}^\top
            \begin{pmatrix}
                2 & 1 & -1\\
                1 & 2 & -1\\
                -1 & -1 & 1
            \end{pmatrix}
            \begin{pmatrix}
                y_1\\
                y_2\\
                x
            \end{pmatrix} + \text{const}\\
            & = -\frac{1}{2} \left( 2 y_1^2 + 2 y_2^2 + x^2 + 2 y_1 y_2 - 2 y_1 x - 2y_2 x \right) + \text{const}
        \end{align}
        
        Vamos come\'car com a equa\'c\~ao para $\bar{p}(y_1|x)$. \'E mais f\'acil trabalhar no dom\'inio logar\'itmico, onde obtemos:
        \begin{align}
            \log  \bar{p}(y_1|x) & =  \E_{q_2(y_2|x)}\left[ \log p(y_1, y_2, x)\right] + \text{const}\\
            & = -\frac{1}{2} \E_{q_2(y_2|x)}\left[  2 y_1^2 + 2 y_2^2 + x^2 + 2 y_1 y_2 - 2 y_1 x - 2y_2 x \right] + \text{const}\\
            & = -\frac{1}{2} \left(2 y_1^2 + 2 y_1 \E_{q_2(y_2|x)}[y_2] - 2 y_1 x \right) + \text{const}\\
            & = -\frac{1}{2} \left( 2 y_1^2 + 2y_1 m_2 - 2 y_1 x \right) + \text{const}\\
            & = -\frac{1}{2} \left( 2 y_1^2 - 2 y_1 (x-m_2) \right) + \text{const}
            \label{eq:barp1-alt}
        \end{align}
        onde absorvemos todos os termos que n\~ao envolvem $y_1$ na constante. Al\'em disso, definimos $\E_{q_2(y_2|x)}[y_2]=m_2$.
        
        Note que uma densidade Gaussiana arbitr\'aria $\Gauss(y; m, \sigma^2)$
        com m\'edia $m$ e vari\^ancia $\sigma^2$ pode ser escrita no
        dom\'inio logar\'itmico como
        \begin{align}
            \log \Gauss(y; m, \sigma^2) & = -\frac{1}{2}\frac{(y-m)^2}{\sigma^2} + \text{const}\\
            & = -\frac{1}{2} \left( \frac{y^2}{\sigma^2} -2y \frac{m}{\sigma^2} \right) + \text{const}
            \label{eq:loggauss-alt}
        \end{align}
        A compara\'c\~ao com \eqref{eq:barp1-alt} mostra que $\bar{p}(y_1|x)$, e
        portanto $q_1(y_1 | x)$, \'e Gaussiana com vari\^ancia e m\'edia iguais a
        \begin{align}
            \sigma_1^2 &= \frac{1}{2} & m_1 & = \frac{1}{2}(x-m_2)
        \end{align}
        Note que n\~ao fizemos uma hip\'otese de Gaussianidade sobre $q_1(y_1 |
        x)$. A $q_1(y_1 | x)$ \'otima acaba sendo Gaussiana porque
        o modelo $p(y_1, y_2, x)$ \'e Gaussiano.
        
        A equa\'c\~ao para $\bar{p}(y_2|x)$ fornece analogamente
        \begin{align}
            \log  \bar{p}(y_2|x) & =  \E_{q_1(y_1|x)}\left[ \log p(y_1, y_2, x)\right] + \text{const}\\
            & = -\frac{1}{2} \E_{q_1(y_1|x)}\left[  2 y_1^2 + 2 y_2^2 + x^2 + 2 y_1 y_2 - 2 y_1 x - 2y_2 x \right] + \text{const}\\
            & = -\frac{1}{2} \left(2 y_2^2 + 2 \E_{q_1(y_1|x)}[y_1] y_2 - 2y_2 x \right) + \text{const}\\
            & = -\frac{1}{2} \left(2 y_2^2 + 2 m_1 y_2 - 2 y_2 x \right) + \text{const}\\
            & = -\frac{1}{2} \left(2 y_2^2 - 2 y_2(x-m_1) \right) + \text{const}
            \label{eq:barp2-alt}
        \end{align}
        onde absorvemos todos os termos que n\~ao envolvem $y_2$ na
        constante. Al\'em disso, definimos $\E_{q_1(y_1|x)}[y_1]=m_1$. Com
        \eqref{eq:loggauss-alt}, isso define uma distribui\'c\~ao Gaussiana com 
        vari\^ancia e m\'edia iguais a
        \begin{align}
            \sigma_2^2 &= \frac{1}{2} & m_2 & = \frac{1}{2}(x-m_1)
        \end{align}
        Portanto, as distribui\'c\~oes variacionais marginais \'otimas $q_1(y_1|x)$ e
        $q_2(y_2|x)$ s\~ao ambas Gaussianas com vari\^ancia igual a $1/2$. Suas
        m\'edias satisfazem
        \begin{align}
            m_1 & = \frac{1}{2}(x-m_2) & m_2  & = \frac{1}{2}(x-m_1)
        \end{align}
        Estas s\~ao duas equa\'c\~oes para duas inc\'ognitas. Podemos resolv\^e-las da seguinte forma   
        \begin{align}
            2 m_1 &= x - m_2 \\
            & = x -\frac{1}{2}(x-m_1)\\
            4 m_1 & = 2 x -x + m_1 \\
            3 m_1 & = x\\
            m_1 & = \frac{1}{3} x
        \end{align}
        Logo
        \begin{equation}
            m_2 = \frac{1}{2} x - \frac{1}{6} x = \frac{2}{6} x = \frac{1}{3} x
        \end{equation}
        Em resumo, encontramos
        \begin{align}
            q_1(y_1 | x) & = \Gauss\left(y_1; \frac{x}{3}, \frac{1}{2}\right) &  q_2(y_2 | x) & = \Gauss\left(y_2; \frac{x}{3}, \frac{1}{2}\right)
        \end{align}
        e a distribui\'c\~ao variacional \'otima $q(y_1, y_2|x) =
        q_1(y_1|x) q_2(y_2|x)$ \'e Gaussiana. Fizemos a hip\'otese de campo m\'edio
        (independ\^encia), mas n\~ao a hip\'otese de Gaussianidade. A Gaussianidade
        da distribui\'c\~ao variacional \'e uma consequ\^encia da Gaussianidade
        do modelo $p(y_1, y_2, x)$.
        
        A compara\'c\~ao com a posterior real mostra que a distribui\'c\~ao variacional
        de campo m\'edio $q(y_1, y_2 | x)$ possui a mesma m\'edia, mas
        ignora a correla\'c\~ao e subestima as vari\^ancias marginais. A posterior real
        e a aproxima\'c\~ao de campo m\'edio s\~ao mostradas na Figura \ref{fig:meanfield-alt}.
        
        \begin{figure}[h!]
            \centering
            \includegraphics[width=0.75\textwidth]{meanfield}
            \caption{\label{fig:meanfield-alt}Em azul: posterior real correlacionada. Em vermelho: aproxima\'c\~ao de campo m\'edio.}
        \end{figure}
    \end{solution}
    
    
    
\end{exenumerate}
% --------------------------------------------------
\ex{Aproxima\'c\~ao variacional da posterior I}

Vimos que maximizar o limite inferior da evid\^encia (ELBO) em
rela\'c\~ao \`a distribui\'c\~ao variacional $q$ minimiza a
diverg\^encia Kullback-Leibler para a posterior real $p$. Aqui assumimos
que $q$ e $p$ s\~ao fun\'c\~oes de densidade de probabilidade, de modo que a
diverg\^encia Kullback-Leibler entre elas \'e definida como
\begin{equation}
    \KL(q || p) = \int q(\x) \log \frac{q(\x)}{p(\x)} \ud \x = \E_q \left[\log \frac{q(\x)}{p(\x)}\right].
\end{equation}



\begin{exenumerate}
    
    \item Voc\^e pode assumir aqui que $\x$ \'e unidimensional, de modo que $p$
    e $q$ sejam densidades univariadas. Considere o caso em que $p$ \'e
    uma densidade bimodal, mas as densidades variacionais $q$ s\~ao
    unimodais. Esbo\'ce uma figura que mostre $p$ e uma distribui\'c\~ao
    variacional $q$ que foi aprendida minimizando $\KL(q || p)$. Explique qualitativamente por que a $q$ esbo\'cada minimiza
    $\KL(q || p)$.
    
    \begin{solution}
        
        Um poss\'ivel esbo\'co \'e mostrado na figura abaixo.
        
        \begin{figure}[h!]
            \centering
            \includegraphics[width = 0.5 \textwidth]{KL-asym}
        \end{figure}
        
        Explica\'c\~ao: Podemos dividir o dom\'inio de $p$ e $q$ nas áreas
        onde $p$ \'e pequena (zero) e naquelas onde $p$ possui massa significativa.
        Como o objetivo apresenta $q$ no numerador enquanto $p$ est\'a
        no denominador, uma $q$ \'otima precisa ser zero onde $p$ \'e
        zero. Caso contr\'ario, incorreria em uma grande penalidade (divis\~ao por
        zero). Como tomamos a esperan\'ca em rela\'c\~ao a $q$, no entanto,
        regi\~oes onde $p>0$ n\~ao precisam ser cobertas por $q$; deix\'a-las
        de fora n\~ao incorre em penalidade. Portanto, as $q$ unimodais \'otimas
        cobrem apenas um pico da bimodal $p$.
        
    \end{solution}
    
    
    \item Assuma que a posterior real $p(\x) = p(x_1, x_2)$ se fatorize em duas Gaussianas de m\'edia zero e vari\^ancias $\sigma_1^2$ e $\sigma_2^2$,
    \begin{equation}
        p(x_1,x_2) = \frac{1}{\sqrt{2\pi \sigma_1^2}} \exp\left[-\frac{x_1^2}{2 \sigma_1^2}\right]\frac{1}{\sqrt{2\pi \sigma_2^2}} \exp\left[-\frac{x_2^2}{2 \sigma_2^2}\right].
    \end{equation}
    Assuma ainda que a densidade variacional $q(x_1,x_2; \lambda^2)$ seja parametrizada como
    \begin{equation}
        q(x_1,x_2;\lambda^2) = \frac{1}{2\pi \lambda^2} \exp\left[-\frac{x_1^2+x_2^2}{2 \lambda^2}\right]
    \end{equation}
    onde $\lambda^2$ \'e o par\^ametro variacional que \'e aprendido
    minimizando $\KL(q || p)$. Se $\sigma^2_2$ for muito
    maior que $\sigma^2_1$, voc\^e espera que $\lambda^2$ esteja mais pr\'oximo de $\sigma_2^2$
    ou de $\sigma_1^2$? Forne\'ca uma explica\'c\~ao.
    
    \begin{solution}
        O par\^ametro variacional aprendido estar\'a mais pr\'oximo de $\sigma_1^2$
        (a menor das duas $\sigma_i^2$).
        
        Explica\'c\~ao: Primeiro, note que as $\sigma_i^2$ s\~ao as vari\^ancias
        ao longo dos dois eixos diferentes, e que $\lambda^2$ \'e a vari\^ancia
        \'unica para ambos $x_1$ e $x_2$. O objetivo penaliza $q$ se ela
        for n\~ao nula onde $p$ \'e pequena (veja acima). O par\^ametro variacional
        $\lambda^2$ ser\'a assim ajustado durante o aprendizado para que a vari\^ancia
        de $q$ esteja pr\'oxima da menor das duas $\sigma_i^2$.
        
    \end{solution}
    
\end{exenumerate}

% --------------------------------------------------

\ex{Aproxima\'c\~ao variacional da posterior II}

Vimos que maximizar o limite inferior da evid\^encia (ELBO) em
rela\'c\~ao \`a distribui\'c\~ao variacional minimiza a
diverg\^encia Kullback-Leibler para a posterior real. Investigamos aqui
a natureza da aproxima\'c\~ao se a fam\'ilia de distribui\'c\~oes
variacionais n\~ao incluir a posterior real.

\begin{exenumerate}
    
    \item Assuma que a posterior real para $\x = (x_1, x_2)$ seja dada por
    \begin{equation}
        p(\x) = \normal(x_1 ; \sigma_1^2)\normal(x_2 ; \sigma_2^2) 
    \end{equation}
    e que nossa distribui\'c\~ao variacional $q(\x; \lambda^2)$ seja
    \begin{equation}
        q(\x; \lambda^2) = \normal(x_1 ; \lambda^2)\normal(x_2 ; \lambda^2),
    \end{equation}
    onde $\lambda > 0$ \'e o par\^ametro variacional. Forne\'ca uma
    equa\'c\~ao para
    \begin{equation}
        J(\lambda) = \KL(q(\x; \lambda^2) || p(\x)),
    \end{equation}
    onde voc\^e pode omitir termos aditivos que n\~ao dependam de $\lambda$. 
    
    \begin{solution}
        
        Escrevemos
        \begin{align}
            \KL(q(\x; \lambda^2) || p(\x)) & = \E_q \left[ \log  \frac{q(\x; \lambda^2)}{p(\x)} \right]\\
            & = \E_q \log q(\x; \lambda^2) - \E_q \log p(\x)\\
            & = \E_q \log  \normal(x_1 ; \lambda^2) + \E_q \log \normal(x_2 ; \lambda^2) \nonumber \\
            & \phantom{=} - \E_q \log \normal(x_1 ; \sigma_1^2) - \E_q \log \normal(x_2 ; \sigma_2^2)
        \end{align}
        
        Temos ainda
        \begin{align}
            \E_q \log  \normal(x_i ; \lambda^2) & = \E_q \log \left[ \frac{1}{\sqrt{2\pi \lambda^2}} \exp \left[-\frac{x_i^2}{2 \lambda^2} \right] \right]\\
            & = \log \left[ \frac{1}{\sqrt{2\pi \lambda^2}}\right] -\E_q \left[\frac{x_i^2}{2 \lambda^2}\right]\\
            & = -\log \lambda - \frac{\lambda^2}{2\lambda^2} + \text{const}\\
            & = -\log \lambda - \frac{1}{2} + \text{const}\\
            & = -\log \lambda + \text{const}
        \end{align}
        onde usamos que, para $x_i$ de m\'edia zero, $\E_q [x_i^2] = \var(x_i) = \lambda^2$. 
        
        Similarmente obtemos
        \begin{align}
            \E_q \log \normal(x_i ; \sigma_i^2) & = \E_q \log \left[ \frac{1}{\sqrt{2\pi \sigma_i^2}} \exp \left[-\frac{x_i^2}{2 \sigma_i^2} \right] \right]\\
            & = \log \left[ \frac{1}{\sqrt{2\pi \sigma_i^2}}\right] -\E_q \left[\frac{x_i^2}{2 \sigma_i^2}\right]\\
            & = -\log \sigma_i - \frac{\lambda^2}{2\sigma_i^2} + \text{const}\\
            & = - \frac{\lambda^2}{2\sigma_i^2} + \text{const}
        \end{align}
        
        Temos assim
        \begin{align}
            \KL(q(\x; \lambda^2 || p(\x))  & =  -2 \log \lambda + \lambda^2\left(\frac{1}{2\sigma_1^2}+\frac{1}{2\sigma_2^2}\right)+\text{const}
        \end{align}     
        
    \end{solution}
    
    
    
    \item Determine o valor de $\lambda$ que minimiza $J(\lambda) =
    \KL(q(\x; \lambda^2) || p(\x))$. Interprete o resultado e
    relacione-o com propriedades da diverg\^encia Kullback-Leibler.
    
    \begin{solution}
        
        Tomando as derivadas de $J(\lambda)$ em rela\'c\~ao a $\lambda$, obtemos
        \begin{align}
            \frac{\partial J(\lambda)}{\partial \lambda} & = -\frac{2}{\lambda} + \lambda \left(\frac{1}{\sigma_1^2}+\frac{1}{\sigma_2^2}\right)
        \end{align}
        Igualando a zero, resulta em
        \begin{align}
            \frac{1}{\lambda^2} &= \frac{1}{2}\left(\frac{1}{\sigma_1^2}+\frac{1}{\sigma_2^2}\right)
        \end{align}
        de modo que
        \begin{align}
            \lambda^2 = 2 \frac{\sigma_1^2 \sigma_2^2}{\sigma_1^2+\sigma_2^2}
        \end{align}
        ou
        \begin{equation}
            \lambda = \sqrt{2} \sqrt{\frac{\sigma_1^2 \sigma_2^2}{\sigma_1^2+\sigma_2^2}}
        \end{equation}
        Este \'e um m\'inimo porque a segunda derivada de $J(\lambda)$
        \begin{equation}
            \frac{\partial^2 J(\lambda)}{\partial \lambda^2} = \frac{2}{\lambda^2} + \left(\frac{1}{\sigma_1^2}+\frac{1}{\sigma_2^2}\right)
        \end{equation}
        \'e positiva para todo $\lambda > 0$.
        
        O resultado possui uma explica\'c\~ao intuitiva: a vari\^ancia \'otima
        $\lambda^2$ \'e a m\'edia harm\^onica das vari\^ancias $\sigma_i^2$ da
        posterior real. Em outras palavras, a precis\~ao \'otima
        $1/\lambda^2$ \'e dada pela m\'edia das precis\~oes
        $1/\sigma_i^2$ das duas dimens\~oes.
        
        Se as vari\^ancias n\~ao forem iguais, por exemplo, se $\sigma_2^2 >
        \sigma_1^2$, vemos que a vari\^ancia \'otima da distribui\'c\~ao variacional
        estabelece um compromisso entre dois tipos de penalidades na diverg\^encia KL: a penalidade de ter um ajuste ruim porque a
        distribui\'c\~ao variacional ao longo da dimens\~ao dois \'e muito estreita; e, ao longo da dimens\~ao um, a penalidade para a distribui\'c\~ao variacional ser diferente de zero quando $p$ \'e pequena. 
        
    \end{solution}
    
\end{exenumerate}


% --------------------------------------------------------------- %
\ex{Algoritmo EM para modelos de mistura}
\label{ex:EM-mixture-models-alt}
Modelos de mistura s\~ao modelos estat\'isticos da forma
\begin{equation}
    p(\x; \thetab) = \sum_{k=1}^K \pi_k p_k(\x; \thetab_k)
    \label{eq:mixture-model-alt}
\end{equation}
onde cada $p_k(\x; \thetab_k)$ \'e em si um modelo estat\'istico parametrizado por
$\thetab_k$ e os $\pi_k \ge 0$ s\~ao pesos de mistura que somam um. Os
par\^ametros $\thetab$ do modelo de mistura consistem nos par\^ametros $\thetab_k$
de cada componente de mistura e nos pesos de mistura $\pi_k$, isto \'e, $\thetab =
(\thetab_1, \ldots, \thetab_K, \pi_1, \ldots, \pi_K)$. Um exemplo \'e uma mistura
de Gaussianas onde cada $p_k(\x; \thetab_k)$ \'e uma Gaussiana com par\^ametros dados
pelo vetor de m\'edia $\mub_k$ e uma matriz de covari\^ancia $\Sigmab_k$.

O modelo de mistura em \eqref{eq:mixture-model-alt} pode ser considerado a
distribui\'c\~ao marginal de um modelo de vari\'avel latente $p(\x, h; \thetab)$ onde $h$
\'e uma vari\'avel n\~ao observada que assume valores $1, \ldots, K$ e $p(h=k) =
\pi_k$. Definindo $p(\x|h=k; \thetab) = p_k(\x; \thetab_k)$, o modelo de vari\'avel
latente correspondente a \eqref{eq:mixture-model-alt} \'e, portanto,
\begin{equation}
    p(\x, h=k; \thetab) = p(\x|h=k; \thetab)p(h=k) = \pi_k p_k(\x; \thetab_k).
    \label{eq:latent-var-model-long-alt}
\end{equation}
Em particular, note que marginalizar $h$ fornece $p(\x; \thetab)$ em
\eqref{eq:mixture-model-alt}.



\begin{exenumerate}
    \item Verifique que o modelo de vari\'avel latente em \eqref{eq:latent-var-model-long-alt} pode ser escrito como
    \begin{equation}
        p(\x, h; \thetab) = \prod_{k=1}^K \left[  \pi_k  p_k(\x; \thetab_k) \right]^{\ind(h=k)}
        \label{eq:latent-var-model-alt}
    \end{equation}
    onde $h$ assume valores em $1, \ldots, K$.
    
    \begin{solution}
        Como $\ind(h=k)$ \'e um se $h=k$ e zero caso contr\'ario, temos
        \begin{equation}
            p(\x, h=j; \thetab) =  \prod_{k=1}^K \left[  \pi_k  p_k(\x; \thetab_k) \right]^{\ind(j=k)} = \pi_j p_j(\x ; \thetab_j)
        \end{equation}
        para qualquer $j \in \{1, \ldots, K\}$, o que coincide com
        \eqref{eq:latent-var-model-long-alt}.
    \end{solution}
    
    \item Como o modelo de mistura em \eqref{eq:mixture-model-alt} pode ser visto como a
    marginal de um modelo de vari\'avel latente, podemos usar o algoritmo de esperan\'ca e maximiza\'c\~ao
    (EM) para estimar os par\^ametros $\thetab$.
    
    Para um modelo geral $p(\data,\h; \thetab)$ onde $\data$ s\~ao os dados observados
    e $\h$ as vari\'aveis n\~ao observadas correspondentes, o algoritmo EM itera entre
    computar a esperan\'ca da log-verossimilhan\'ca de dados completos $J^l(\thetab)$ e maximiz\'a-la em rela\'c\~ao a $\thetab$: 
    \begin{align}
        \textbf{Passo E na itera\'c\~ao l:}\quad J^l(\thetab) &= \E_{p(\h | \data; \thetab^l)} [ \log p(\data,\h; \thetab)]\\
        \textbf{Passo M na itera\'c\~ao l:}\quad \thetab^{l+1} &= \argmax_{\thetab} J^l(\thetab) 
    \end{align}
    Aqui $\thetab^l$ \'e o valor de $\thetab$ na $l$-\'esima itera\'c\~ao. Ao
    resolver o problema de otimiza\'c\~ao, tamb\'em precisamos levar em conta as
    restri\'c\~oes nos par\^ametros, por exemplo, que os $\pi_k$ correspondam a uma fmp.
    
    Assuma que os dados $\data$ consistam em $n$ pontos de dados iid $\x_i$,
    que cada $\x_i$ possua associada a si uma vari\'avel n\~ao observada escalar $h_i$,
    e que as tuplas $(\x_i, h_i)$ sejam todas iid. O que \'e $J^l(\thetab)$ sob
    estas hip\'oteses adicionais?
    
    \begin{solution}
        Como as $(\x_i, h_i)$ s\~ao iid, temos que $p(\data,\h; \thetab) =
        \prod_{i=1}^n p(\x_i, h_i; \thetab)$. Portanto,
        \begin{align}
            J^l(\thetab) &= \E_{p(\h | \data; \thetab^l)} [ \log p(\data,\h; \thetab)]\\
            & = \E_{p(\h | \data; \thetab^l)}\left[ \sum_{i=1}^n \log  p(\x_i, h_i; \thetab)\right]\\
            & = \sum_{i=1}^n \E_{p(\h | \data; \thetab^l)}[ \log  p(\x_i, h_i; \thetab)]\\
            & =  \sum_{i=1}^n \E_{p(h_i | \data; \thetab^l)}[ \log  p(\x_i, h_i; \thetab)]\\
            & =  \sum_{i=1}^n \E_{p(h_i | \x_i; \thetab^l)}[ \log  p(\x_i, h_i; \thetab)] \label{eq:Jl-intermediate-alt}
        \end{align}
        onde no pen\'ultimo passo usamos que cada $\log p(\x_i, h_i;
        \thetab)]$ envolve apenas uma vari\'avel latente $h_i$, de modo que s\'o precisamos
        tomar a esperan\'ca em rela\'c\~ao a $p(h_i | \data; \thetab^l)$, e no \'ultimo passo,
        usamos que $h_i \independent \x_j$ para $j \neq i$.
        
    \end{solution}
    
    \item Mostre que para o modelo de vari\'avel latente em \eqref{eq:latent-var-model-alt},
    $J^l(\thetab)$ \'e igual a
    \begin{align}
        J^l(\thetab) & = \sum_{i=1}^n  \sum_{k=1}^K w_{ik}^l  \log[  \pi_k  p_k(\x_i; \thetab_k) ],\label{eq:Jl-expression-alt}\\
        w_{ik}^l & =  \frac{ \pi^l_k p_k(\x_i; \thetab^l_k)}{ \sum_{k=1}^K \pi^l_k p_k(\x_i; \thetab^l_k)}\label{eq:w-def-alt}
    \end{align}
    Note que os $w_{ik}^l$ s\~ao definidos em termos dos par\^ametros $\pi_k^l$ e
    $\thetab_k^l$ da itera\'c\~ao $l$. Eles s\~ao iguais \`as probabilidades condicionais
    $p(h=k|\x_i; \thetab^l)$, isto \'e, a probabilidade de que $\x_i$
    tenha sido amostrada da componente $p_k(\x_i; \thetab^l_k)$. 
    
    \begin{solution}
        Consideramos um \'unico termo $ \E_{p(h| \x; \thetab^l)}[ \log p(\x, h;
        \thetab)]$ em \eqref{eq:Jl-intermediate-alt}.
        
        Dada a forma do modelo em \eqref{eq:latent-var-model-alt}, temos que
        \begin{align}
            \log p(\x, h; \thetab) = \sum_{k=1}^K \ind(h=k) \log[  \pi_k  p_k(\x; \thetab_k) ]
        \end{align}
        e portanto
        \begin{align}
            \E_{p(h| \x; \thetab^l)}[ \log p(\x, h; \thetab)] &=  \E_{p(h| \x; \thetab^l)}\left[ \sum_{k=1}^K \ind(h=k) \log[  \pi_k  p_k(\x; \thetab_k) ] \right]\\
            &=  \sum_{k=1}^K  \E_{p(h| \x; \thetab^l)}\left[ \ind(h=k) \right] \log[  \pi_k  p_k(\x; \thetab_k) ]\\
            & = \sum_{k=1}^K p(h=k|\x; \thetab^l)  \log[  \pi_k  p_k(\x; \thetab_k) ]
        \end{align}
        onde usamos que a esperan\'ca sobre um evento indicadora \'e igual \`a
        probabilidade do evento ocorrer, isto \'e, $\E_{p(h| \x; \thetab^l)}\left[
        \ind(h=k) \right] =  p(h=k|\x; \thetab^l)$.
        
        A probabilidade $p(h=k|\x; \thetab^l)$ pode ser determinada via a regra do produto
        (regra de Bayes) e as Equa\'c\~oes \eqref{eq:latent-var-model-long-alt} e \eqref{eq:mixture-model-alt}
        \begin{align}
            p(h=k|\x; \thetab^l) &= \frac{p(\x, h=k, \thetab^l)}{p(\x; \thetab^l)}\\
            & = \frac{ \pi^l_k p_k(\x; \thetab^l_k)}{ \sum_{k=1}^K \pi^l_k p_k(\x; \thetab^l_k)}
            \label{eq:posterior-def-alt}
        \end{align}
        Note que o sobrescrito $^l$ indica que os $\pi_k^l$ s\~ao os pesos de mistura
        e os $\thetab^l_k$ os par\^ametros do modelo da itera\'c\~ao $l$.
        
        O objetivo $J^l(\thetab)$ soma sobre $n$ termos $ \E_{p(h| \x_i; \thetab^l)}[ \log p(\x_i, h;
        \thetab)]$. Denotemos $p(h=k|\x_i; \thetab^l)$ de \eqref{eq:posterior-def-alt} por $w_{ik}^l$ de modo que
        \begin{equation}
            \E_{p(h| \x_i; \thetab^l)}[ \log p(\x_i, h; \thetab)] =  \sum_{k=1}^K w_{ik}^l  \log[  \pi_k  p_k(\x; \thetab_k) ]
        \end{equation}
        e
        \begin{align}
            J^l(\thetab) & = \sum_{i=1}^n  \sum_{k=1}^K w_{ik}^l  \log[  \pi_k  p_k(\x_i; \thetab_k) ].
        \end{align}
        O objetivo $J^l(\thetab)$ assume a forma de uma log-verossimilhan\'ca ponderada. Em
        mais detalhes, como $\sum_k w_{ik}^l=1$ para todos os pontos de dados $\x_i$ (e
        $w_{ik}^l \ge 0 $), $ \sum_{k=1}^K w_{ik}^l \log[ \pi_k p_k(\x_i; \thetab_k)
        ]$ \'e uma combina\'c\~ao convexa. Isso significa que os diferentes componentes do
        modelo de mistura competem entre si: pesos maiores para algumas
        componentes significam pesos menores para outras. No caso extremo, algumas
        componentes podem contribuir de forma negligenci\'avel para o $i$-\'esimo termo da
        log-verossimilhan\'ca.
        
        Os pesos $w_{ik}^l$ s\~ao s\^o vezes, em particular para misturas de
        Gaussianas, chamados de ``atribui\'c\~oes suaves'' (*soft-assignments*) porque especificam em que medida
        um ponto de dado $\x_i$ ``pertence'' a uma componente de mistura $p_k$.
        Alternativamente, podemos interpretar os $w_{ik}^l$ como as
        ``responsabilidades'' de cada componente de mistura $p_k$ por um ponto de dado $\x_i$.
        
        Em alguns casos, e.g.\ por raz\~oes computacionais, podemos determinar qual dos
        $K$ pesos $w_{i1}^l, \ldots, w_{iK}^l$ \'e o maior e ent\~ao defini-lo como
        um, enquanto definimos os outros pesos como zero. Isso corresponde a
        ``atribui\'c\~oes r\'igidas'' (*hard-assignments* e ``hard EM'') onde um ponto de dado $\x_i$ \'e
        exclusivamente atribu\'ido a uma \'unica componente de mistura $p_k$. 
    \end{solution}
    
    \item Assuma que as diferentes componentes de mistura $p_k(\x; \thetab_k), k=1,
    \ldots, K$ n\~ao compartilhem nenhum par\^ametro. Mostre que os valores de par\^ametros atualizados
    $\thetab_k^{l+1}$ s\~ao dados por estimativas de m\'axima verossimilhan\'ca ponderadas.
    
    \begin{solution}
        Intercambiamos a ordem dos somat\'orios em \eqref{eq:Jl-expression-alt} de modo
        que
        \begin{align}
            J^l(\thetab) & = \sum_{k=1}^K \sum_{i=1}^n  w_{ik}^l  \log[  \pi_k  p_k(\x_i; \thetab_k) ]\\
            & =  \sum_{k=1}^K \sum_{i=1}^n  w_{ik}^l \log \pi_k + \sum_{k=1}^K \underbrace{\sum_{i=1}^n  w_{ik}^l \log p_k(\x_i; \thetab_k)}_{\ell_k^l(\thetab_k)}
            \label{eq:Jl-interchanged-alt}
        \end{align}
        Quando atualizamos os par\^ametros $\thetab_k$ das componentes de mistura, o
        primeiro termo \'e uma constante. O segundo termo \'e uma soma sobre as
        log-verossimilhan\'cas ponderadas $\ell_k^l(\thetab_k)$, uma para cada componente de mistura. Se as
        componentes de mistura n\~ao compartilham par\^ametros, temos ent\~ao
        \begin{equation}
            \thetab_k^{l+1} =  \argmax_{\thetab_k}  J^l(\thetab) = \argmax_{\thetab_k} \ell_k^l(\thetab_k)
        \end{equation}
        Isso significa que podemos calcular $\thetab_k^{l+1}$ como se realiz\'assemos a estimativa de
        m\'axima verossimilhan\'ca para o modelo $p_k(\x; \thetab_k)$, exceto pelo fato de
        que os pontos de dados $\x_i$ s\~ao ponderados pelos $w_{ik}^l$.
    \end{solution}
    
    \item Mostre que maximizar $J^l(\thetab)$ em rela\'c\~ao aos pesos de mistura
    $\pi_k$ fornece a regra de atualiza\'c\~ao
    \begin{align}
        \pi_k^{l+1} &=  \frac{1}{n}\sum_{i=1}^n w_{ik}^l
    \end{align}
    \begin{solution}
        Come\'camos com \eqref{eq:Jl-expression-alt} e descartamos os termos aditivos que n\~ao
        dependem de $\pi_k$. Como
        \begin{align}
            J^l(\thetab) & =  \sum_{i=1}^n  \sum_{k=1}^K w_{ik}^l  \log \pi_k + \text{termos que n\~ao dependem dos } \pi_k
        \end{align}
        podemos focar no objetivo
        \begin{align}
            J^l_\pi(\pi_1, \ldots, \pi_K) & = \sum_{i=1}^n  \sum_{k=1}^K w_{ik}^l  \log \pi_k\\ 
            &= \sum_{k=1}^K \underbrace{\left( \sum_{i=1}^n w_{ik}^l \right)}_{\omega_k^l}  \log \pi_k \\
            & =  \sum_{k=1}^K \omega_k^l \log \pi_k.
        \end{align}
        Levando em conta que os $\pi_k = p(h=k)$ definem uma fmp, o problema de otimiza\'c\~ao a ser resolvido \'e
        \begin{align}
            \text{maximizar} \quad&  \sum_{k=1}^K \omega_k^l \log \pi_k\\
            \text{sujeito a}\quad &  \pi_k \ge 0 \\
            &  \sum_{k=1}^K \pi_k = 1 
        \end{align}
        O problema de otimiza\'c\~ao restrito poderia ser resolvido via multiplicadores de
        Lagrange. Mas aqui adotamos outra abordagem e resolvemos o problema de otimiza\'c\~ao
        expressando-o em termos de um problema de minimiza\'c\~ao da diverg\^encia KL.
        
        Primeiro, note que os $\pi_k$ que maximizam $J^l_\pi(\pi_1, \ldots,
        \pi_K)$ tamb\'em maximizar\~ao o objetivo reescalonado
        \begin{align}
            \frac{1}{\sum_{k=1}^K \omega_k^l}  J_\pi^l(\pi_1, \ldots, \pi_K) & =  \frac{1}{\sum_{k=1}^K \omega_k^l}\sum_{k=1}^K \omega_k^l \log \pi_k\\
            &=   \sum_{k=1}^K q_k^l \log \pi_k
        \end{align}
        onde introduzimos
        \begin{equation}
            q_k^l = \frac{\omega_k^l}{\sum_{k=1}^K \omega_{k}^l}.
        \end{equation}
        Os $q_k^l$ s\~ao n\~ao-negativos e somam um. Portanto, podemos considerar que eles
        definem uma fmp.
        
        Segundo, note que os $\pi_k$ que maximizam $J_\pi^l(\pi_1, \ldots,
        \pi_K)$ tamb\'em maximizar\~ao
        \begin{align}
            \sum_{k=1}^K q_k^l \log \pi_k - \sum_{k=1}^K q_k^l \log q_k^l & = \sum_{k=1}^K q_k^l \log \frac{\pi_k}{q_k^l}\\
            & = -\sum_{k=1}^K q_k^l \log \frac{q_k^l}{\pi_k}\\
            & = -\KL(q^l, \pi)
        \end{align}
        j\'a que adicionar constantes n\~ao altera a solu\'c\~ao. Assim, os
        $\pi_k$ \'otimos s\~ao dados pela fmp $\pi$ que minimiza a diverg\^encia KL
        $\KL(q^l, \pi)$. Isso significa que os $\pi_k$ \'otimos s\~ao
        \begin{align}
            \pi_k & = q_k^l= \frac{\omega_k^l}{\sum_{k=1}^K \omega_{k}^l} = \frac{\sum_{i=1}^n w_{ik}^l}{\sum_{k=1}^K \sum_{i=1}^n w_{ik}^l }.
        \end{align}
        O denominador pode ser simplificado notando que, com \eqref{eq:w-def-alt},
        $\sum_{k=1}^K w_{ik}^l=1$, de modo que
        \begin{align}
            \sum_{k=1}^K \sum_{i=1}^n w_{ik}^l = \sum_{i=1}^n \sum_{k=1}^K w_{ik}^l = n
        \end{align}   
        A regra de atualiza\'c\~ao solicitada \'e portanto
        \begin{align}
            \pi_k^{l+1} &=  \frac{1}{n}\sum_{i=1}^n w_{ik}^l
        \end{align}
        A regra de atualiza\'c\~ao n\~ao depende diretamente do modelo estat\'istico $p_k(\x;
        \thetab_k)$ que podemos escolher para as componentes de mistura. Sua influ\^encia
        ocorre indiretamente atrav\'es dos $w_{ik}^l$.
    \end{solution}
    
    \item \label{ex:EM-mixture-models-summary-alt} Resuma o algoritmo EM para aprender os par\^ametros $\thetab$ do
    modelo de mistura em \eqref{eq:mixture-model-alt} a partir de dados iid $\x_1, \ldots, \x_n$.
    
    \begin{solution}
        Coletamos e resumimos os resultados das quest\~oes anteriores:
        \begin{itemize}
            \item \textbf{Passo E na itera\'c\~ao l:} Compute as probabilidades posteriores (atribui\'c\~oes suaves)
            \begin{align}
                w_{ik}^l & =  \frac{ \pi^l_k p_k(\x_i; \thetab^l_k)}{ \sum_{k=1}^K \pi^l_k p_k(\x_i; \thetab^l_k)} 
            \end{align}
            para todos os pontos de dados $\x_i$ e componentes de mistura $k$. Em seguida, formule
            a fun\'c\~ao objetivo $J^l(\thetab)$
            \begin{align}
                J^l(\thetab) & = \sum_{i=1}^n  \sum_{k=1}^K w_{ik}^l  \log[  \pi_k  p_k(\x_i; \thetab_k) ]
            \end{align}
            \item  \textbf{Passo M na itera\'c\~ao l:} Compute os novos pesos de mistura
            \begin{align}
                \pi_k^{l+1} &=  \frac{1}{n}\sum_{i=1}^n w_{ik}^l
            \end{align}
            Para computar os novos par\^ametros de mistura $\thetab_k^{l+1}$, maximize
            $J^l(\thetab)$ se alguns par\^ametros forem compartilhados ou vinculados. Se os $p_k(\x;
            \thetab_k)$ n\~ao compartilharem par\^ametros, os novos par\^ametros $\thetab_k^{l+1}$
            s\~ao obtidos maximizando uma log-verossimilhan\'ca ponderada para cada componente de
            mistura separadamente:
            \begin{align}
                \thetab_k^{l+1} &= \argmax_{\thetab_k} \sum_{i=1}^n  w_{ik}^l \log p_k(\x_i; \thetab_k)
            \end{align}
            para $k=1, \ldots, K$.
        \end{itemize}
    \end{solution}
    
\end{exenumerate}


% --------------------------------------------------------------- %

\ex{Algoritmo EM para mistura de Gaussianas}
\label{ex:EM-MoG-alt}

Usamos aqui os resultados do \exref{ex:EM-mixture-models-alt} para derivar as regras de atualiza\'c\~ao EM para uma mistura de Gaussianas. Este \'e um modelo de mistura onde cada
componente de mistura \'e uma distribui\'c\~ao Gaussiana, isto \'e,
\begin{align}
    p(\x; \thetab) = \sum_{i=1}^K \pi_k \normal(\x; \mub_k, \Sigmab_k).
\end{align}
Consideramos o caso em que cada $\mub_k$ e $\Sigmab_k$ pode ser individualmente
alterado (sem vincula\'c\~ao de par\^ametros). Os par\^ametros globais do modelo s\~ao dados
por $\mub_k, \Sigmab_k$ e os pesos de mistura $\pi_k \ge 0$, $k=1, \ldots,
K$. Como no caso de modelos de mistura gerais, os pesos de mistura somam um.

\begin{exenumerate}
    
    \item Determine as estimativas de m\'axima verossimilhan\'ca para uma Gaussiana multivariada
    $\normal(\x; \mub, \Sigmab)$ para dados iid $\data = (\x_1, \ldots, \x_n)$ quando
    cada ponto de dado $\x_i$ possui um peso $w_i$. Os pesos s\~ao n\~ao-negativos, mas n\~ao
    somam necessariamente um.
    
    \begin{solution}
        A log-verossimilhan\'ca ponderada \'e
        \begin{align}
            \ell(\mub, \Sigmab) & = \sum_{i=1}^n w_i \log \normal(\x_i; \mub, \Sigmab)\\
            & =  \sum_{i=1}^n w_i \log |\det(2 \pi \Sigmab)|^{-1/2} - \frac{1}{2} \sum_{i=1}^n w_i (\x_i-\mub)^\top\Sigmab^{-1}(\x_i-\mub)
        \end{align}
        Introduzindo os pesos normalizados $W_i = w_i/ \sum_{i=1}^n w_i$, temos
        \begin{align}
            \frac{1}{ \sum_{i=1}^n w_i} \ell(\mub, \Sigmab)  &=  \log |\det(2 \pi \Sigmab)|^{-1/2} - \frac{1}{2} \sum_{i=1}^n W_i (\x_i-\mub)^\top\Sigmab^{-1}(\x_i-\mub)
        \end{align}
        Vamos escrever o termo quadr\'atico
        \begin{align} 
            (\x_i-\mub)^T\Sigmab^{-1}(\x_i-\mub) &= \x_i^\top\Sigmab^{-1}\x_i- 2\x_i^\top \Sigmab^{-1} \mub + \mub^\top \Sigmab^{-1} \mub
        \end{align}
        Logo
        \begin{align}
            \sum_{i=1}^n W_i (\x_i-\mub)^T\Sigmab^{-1}(\x_i-\mub) &= \sum_{i=1}^n W_i \x_i^\top\Sigmab^{-1}\x_i -  2 \sum_{i=1}^n W_i\x_i^\top \Sigmab^{-1} \mub +
            \underbrace{\sum_{i=1}^n W_i}_{=1}    \mub^\top \Sigmab^{-1} \mub  \\
            & = \tr\left[ \left( \sum_{i=1}^n W_i \x_i \x_i^\top \right) \Sigmab^{-1} \right]- 2 \left( \sum_{i=1}^n W_i\x_i  \right)^\top \Sigmab^{-1} \mub +  \mub^\top \Sigmab^{-1} \mub \\
            &= \tr \left( \R \Sigmab^{-1}\right) -2 \b^\top \Sigmab^{-1}\mub +  \mub^\top\Sigmab^{-1} \mub
        \end{align}
        onde $\R = \sum_{i=1}^n W_i \x_i \x_i^\top $ e $\b =  \sum_{i=1}^n
        W_i\x_i$. Portanto,
        \begin{align}
            \frac{1}{ \sum_{i=1}^n w_i} \ell(\mub, \Sigmab)  &=  \log |\det(2 \pi \Sigmab)|^{-1/2} - \frac{1}{2} \tr \left( \R \Sigmab^{-1}\right) + \b^\top \Sigmab^{-1}\mub - \frac{1}{2}  \mub^\top\Sigmab^{-1} \mub
        \end{align}
        Isso possui exatamente a mesma forma que a fun\'c\~ao de verossimilhan\'ca n\~ao ponderada, apenas as
        estat\'isticas suficientes $\R$ e $\b$ s\~ao computadas usando os pesos. Assim,
        as estimativas de m\'axima verossimilhan\'ca, quando expressas em termos de $\R$ e $\b$,
        permanecem as mesmas do caso n\~ao ponderado:
        \begin{align}
            \hat{\mub} & = \b  =  \sum_{i=1}^n W_i\x_i\\
            \hat{\Sigmab} & = \R - \b\b^\top=  \sum_{i=1}^n W_i \x_i \x_i^\top - \b\b^\top
        \end{align}
        Al\'em disso, como
        \begin{align}
            \sum_{i=1}^n W_i (\x_i-\b)(\x_i-\b)^\top  & =  \sum_{i=1}^n W_i \x_i \x_i^\top - \underbrace{\sum_{i=1}^n W_i \x_i}_{\b} \b^\top - \b \underbrace{\sum_{i=1}^n W_i \x_i^\top}_{\b^\top} + \b\b^\top\\
            &=\R - \b \b^\top - \b \b^\top + \b\b^\top\\
            & =  \R - \b\b^\top
        \end{align}
        descobrimos que as estimativas de m\'axima verossimilhan\'ca ponderadas s\~ao a m\'edia ponderada
        e a matriz de covari\^ancia ponderada:
        \begin{align}
            \hat{\mub} & =   \sum_{i=1}^n W_i\x_i&  \hat{\Sigmab} &=    \sum_{i=1}^n W_i (\x_i-\hat{\mub})(\x_i-\hat{\mub})^\top & W_i&= \frac{w_i}{\sum_{i=1}^n w_i}
        \end{align}
    \end{solution}
    
    \item Use os resultados do \exref{ex:EM-mixture-models-alt} para derivar as regras de atualiza\'c\~ao EM para os par\^ametros do modelo de mistura Gaussiana.
    
    \begin{solution}
        Da solu\'c\~ao do \exref{ex:EM-mixture-models-alt}\ref{ex:EM-mixture-models-summary-alt} e das solu\'c\~oes MLE ponderadas
        derivadas, encontramos:
        \begin{itemize}
            \item \textbf{Passo E na itera\'c\~ao l:} Compute as probabilidades posteriores (atribui\'c\~oes suaves)
            \begin{align}
                w_{ik}^l & =  \frac{ \pi^l_k \normal(\x_i; \mub^l_k, \Sigmab_k^l)}{ \sum_{k=1}^K \pi^l_k \normal(\x_i; \mub^l_k, \Sigmab_k^l) } 
            \end{align}
            para todos os pontos de dados $\x_i$ e componentes de mistura $k$.
            \item  \textbf{Passo M na itera\'c\~ao l:}
            \begin{itemize}
                \item Determine as MLEs ponderadas
                \begin{align}
                    \mub_k^{l+1} & =   \sum_{i=1}^n W^l_{ik}\x_i&  \Sigmab_k^{l+1} &=    \sum_{i=1}^n W^l_{ik} (\x_i - \mub_k^{l+1})(\x_i - \mub_k^{l+1})^\top
                \end{align}
                onde $W^l_{ik} = w_{ik}^l/(\sum_{i=1}^n w_{ik}^l)$.
                \item Compute os novos pesos de mistura
                \begin{align}
                    \pi_k^{l+1} &=  \frac{1}{n}\sum_{i=1}^n w_{ik}^l
                \end{align}
            \end{itemize}
        \end{itemize}
        
    \end{solution}
    
\end{exenumerate}